% Appendix A: Reproducibility Setup
%
% Generated 2026-09-08 (Round-9 / session continuation) to fill the
% missing appendix that was flagged by the Tier-6 audit 2026-09-07
% as part of the "Emergency stop" build issue.
%
% Cross-referenced from thesis/main.tex % Appendix A: Reproducibility Setup
%
% Generated 2026-09-08 (Round-9 / session continuation) to fill the
% missing appendix that was flagged by the Tier-6 audit 2026-09-07
% as part of the "Emergency stop" build issue.
%
% Cross-referenced from thesis/main.tex % Appendix A: Reproducibility Setup
%
% Generated 2026-09-08 (Round-9 / session continuation) to fill the
% missing appendix that was flagged by the Tier-6 audit 2026-09-07
% as part of the "Emergency stop" build issue.
%
% Cross-referenced from thesis/main.tex % Appendix A: Reproducibility Setup
%
% Generated 2026-09-08 (Round-9 / session continuation) to fill the
% missing appendix that was flagged by the Tier-6 audit 2026-09-07
% as part of the "Emergency stop" build issue.
%
% Cross-referenced from thesis/main.tex \input{chapters/appendix_a_setup}.
% Wordcount target: ~500 words of substantive prose + commands.

\chapter{Reproducibility Setup}
\label{ch:appendix:setup}

This appendix documents the complete reproducibility chain for the
six-paper thesis substrate: how to clone the repository, install the
Python dependencies, verify the installation, and run the full
verification suite. Every result in the thesis body is reproducible
from the public GitHub release by following the steps below; no
proprietary data or hidden state is required.

\section{Prerequisites}

The repository targets Python 3.10+ and depends on the
\texttt{uv} package manager for reproducible environment installation.
System-level dependencies (GDAL, GEOS, PROJ) are required for the
\texttt{rasterio} and \texttt{geopandas} imports used in the data
loaders; on Debian/Ubuntu these can be installed via the
\texttt{libgdal-dev} / \texttt{libgeos-dev} / \texttt{libproj-dev}
packages, or via the bundled multi-stage \texttt{Dockerfile}
which provisions them automatically.

\section{Installation}

The standard install path uses the project Makefile:

\begin{verbatim}
git clone https://github.com/IvanWeissVanDerPol/satellite-paraguay
cd satellite-paraguay
make install          # uv sync --all-extras + pip install -e .
make bootstrap        # install + data catalog + sample download
\end{verbatim}

The \texttt{bootstrap} target creates the \texttt{.venv/} virtual
environment, installs all Python dependencies pinned in
\texttt{uv.lock}, and downloads the sample Sentinel-2 + Hansen tiles
needed for the smoke tests. A full data bootstrap (real Hansen
v1.11 coverage of Paraguay + Sentinel-2 batch download) is the
\texttt{make data-sentinel} + \texttt{make data-mapbiomas} +
\texttt{make data-hansen} target chain.

\section{Verification}

Three verification commands establish that the installed environment
matches the canonical state of the repository:

\begin{itemize}
    \item \texttt{bash scripts/run\_tests.sh} — full pytest suite
      (~759 tests, ~80\,s wall-clock). Catches regressions in
      \texttt{src/}, \texttt{tests/}, \texttt{scripts/}, and the
      per-paper \texttt{paper.tex} / \texttt{references.bib} files.
    \item \texttt{bash scripts/run\_tests\_fast.sh} — fast inner-loop
      subset (7 test files, ~5\,s). Recommended for post-edit
      sanity checks during development.
    \item \texttt{bash scripts/defense\_check.py} — eight-check defense
      gate that bundles citation resolution, claim integrity,
      LaTeX syntax, ethics gates, and CrossRef DOI verification.
      This is the canonical pre-submission command.
\end{itemize}

All three commands are designed to exit non-zero on any failure; the
\texttt{run\_tests.sh} wrapper uses \texttt{set -euo pipefail} so a
single broken test blocks the whole run.

\section{Data Version Control}

Datasets larger than a few megabytes are tracked via DVC
(\texttt{dvc.yaml}) rather than committed to git. The canonical data
storage layer is the companion repository
\href{https://github.com/IvanWeissVanDerPol/paraguay-geodata-vlm}{\texttt{paraguay-geodata-vlm}}
(local path \texttt{/opt/data/thesis-active/}), which contains the
9-dataset manifest, OSM Paraguay downloader, and SAM +
GroundingDINO annotation pipeline. Per-paper \texttt{DATA\_MANIFEST.md}
files record which datasets are required for each paper; the
\texttt{scripts/validate\_data.py} worker audits whether the local
files actually exist or are missing/synthetic placeholders.

\section{Containerized Reproducibility}

For environments without \texttt{uv} or with restricted network
access, the bundled \texttt{Dockerfile} provides a fully
self-contained build:

\begin{verbatim}
docker build -t satellite-paraguay .
docker run -it --rm \
    -v $(pwd)/data:/app/data \
    satellite-paraguay bash scripts/run_tests.sh
\end{verbatim}

The image is multi-stage (Python 3.10-slim + GDAL system deps +
\texttt{uv}-installed Python deps) and produces a deterministic
environment that matches the \texttt{uv.lock} lock file. A
\texttt{docker-compose.production.yml} file is included for
deploying the FastAPI + Streamlit demo to a remote host.

\section{Continuous Integration}

The repository includes a GitHub Actions workflow
(\texttt{.github/workflows/ci.yml}) that runs the fast pytest subset
+ lint chain (black + flake8 + isort + mypy strict) on every push
and PR. A separate \texttt{latex.yml} workflow compiles each of the
six \texttt{paper.tex} files with TeX Live, ensuring that no paper
silently loses its ability to compile. The CI workflow SHAs were
pinned to v4.4.0 in commit \texttt{0b66057} after the v7.0.1
references were force-pushed upstream.
.
% Wordcount target: ~500 words of substantive prose + commands.

\chapter{Reproducibility Setup}
\label{ch:appendix:setup}

This appendix documents the complete reproducibility chain for the
six-paper thesis substrate: how to clone the repository, install the
Python dependencies, verify the installation, and run the full
verification suite. Every result in the thesis body is reproducible
from the public GitHub release by following the steps below; no
proprietary data or hidden state is required.

\section{Prerequisites}

The repository targets Python 3.10+ and depends on the
\texttt{uv} package manager for reproducible environment installation.
System-level dependencies (GDAL, GEOS, PROJ) are required for the
\texttt{rasterio} and \texttt{geopandas} imports used in the data
loaders; on Debian/Ubuntu these can be installed via the
\texttt{libgdal-dev} / \texttt{libgeos-dev} / \texttt{libproj-dev}
packages, or via the bundled multi-stage \texttt{Dockerfile}
which provisions them automatically.

\section{Installation}

The standard install path uses the project Makefile:

\begin{verbatim}
git clone https://github.com/IvanWeissVanDerPol/satellite-paraguay
cd satellite-paraguay
make install          # uv sync --all-extras + pip install -e .
make bootstrap        # install + data catalog + sample download
\end{verbatim}

The \texttt{bootstrap} target creates the \texttt{.venv/} virtual
environment, installs all Python dependencies pinned in
\texttt{uv.lock}, and downloads the sample Sentinel-2 + Hansen tiles
needed for the smoke tests. A full data bootstrap (real Hansen
v1.11 coverage of Paraguay + Sentinel-2 batch download) is the
\texttt{make data-sentinel} + \texttt{make data-mapbiomas} +
\texttt{make data-hansen} target chain.

\section{Verification}

Three verification commands establish that the installed environment
matches the canonical state of the repository:

\begin{itemize}
    \item \texttt{bash scripts/run\_tests.sh} — full pytest suite
      (~759 tests, ~80\,s wall-clock). Catches regressions in
      \texttt{src/}, \texttt{tests/}, \texttt{scripts/}, and the
      per-paper \texttt{paper.tex} / \texttt{references.bib} files.
    \item \texttt{bash scripts/run\_tests\_fast.sh} — fast inner-loop
      subset (7 test files, ~5\,s). Recommended for post-edit
      sanity checks during development.
    \item \texttt{bash scripts/defense\_check.py} — eight-check defense
      gate that bundles citation resolution, claim integrity,
      LaTeX syntax, ethics gates, and CrossRef DOI verification.
      This is the canonical pre-submission command.
\end{itemize}

All three commands are designed to exit non-zero on any failure; the
\texttt{run\_tests.sh} wrapper uses \texttt{set -euo pipefail} so a
single broken test blocks the whole run.

\section{Data Version Control}

Datasets larger than a few megabytes are tracked via DVC
(\texttt{dvc.yaml}) rather than committed to git. The canonical data
storage layer is the companion repository
\href{https://github.com/IvanWeissVanDerPol/paraguay-geodata-vlm}{\texttt{paraguay-geodata-vlm}}
(local path \texttt{/opt/data/thesis-active/}), which contains the
9-dataset manifest, OSM Paraguay downloader, and SAM +
GroundingDINO annotation pipeline. Per-paper \texttt{DATA\_MANIFEST.md}
files record which datasets are required for each paper; the
\texttt{scripts/validate\_data.py} worker audits whether the local
files actually exist or are missing/synthetic placeholders.

\section{Containerized Reproducibility}

For environments without \texttt{uv} or with restricted network
access, the bundled \texttt{Dockerfile} provides a fully
self-contained build:

\begin{verbatim}
docker build -t satellite-paraguay .
docker run -it --rm \
    -v $(pwd)/data:/app/data \
    satellite-paraguay bash scripts/run_tests.sh
\end{verbatim}

The image is multi-stage (Python 3.10-slim + GDAL system deps +
\texttt{uv}-installed Python deps) and produces a deterministic
environment that matches the \texttt{uv.lock} lock file. A
\texttt{docker-compose.production.yml} file is included for
deploying the FastAPI + Streamlit demo to a remote host.

\section{Continuous Integration}

The repository includes a GitHub Actions workflow
(\texttt{.github/workflows/ci.yml}) that runs the fast pytest subset
+ lint chain (black + flake8 + isort + mypy strict) on every push
and PR. A separate \texttt{latex.yml} workflow compiles each of the
six \texttt{paper.tex} files with TeX Live, ensuring that no paper
silently loses its ability to compile. The CI workflow SHAs were
pinned to v4.4.0 in commit \texttt{0b66057} after the v7.0.1
references were force-pushed upstream.
.
% Wordcount target: ~500 words of substantive prose + commands.

\chapter{Reproducibility Setup}
\label{ch:appendix:setup}

This appendix documents the complete reproducibility chain for the
six-paper thesis substrate: how to clone the repository, install the
Python dependencies, verify the installation, and run the full
verification suite. Every result in the thesis body is reproducible
from the public GitHub release by following the steps below; no
proprietary data or hidden state is required.

\section{Prerequisites}

The repository targets Python 3.10+ and depends on the
\texttt{uv} package manager for reproducible environment installation.
System-level dependencies (GDAL, GEOS, PROJ) are required for the
\texttt{rasterio} and \texttt{geopandas} imports used in the data
loaders; on Debian/Ubuntu these can be installed via the
\texttt{libgdal-dev} / \texttt{libgeos-dev} / \texttt{libproj-dev}
packages, or via the bundled multi-stage \texttt{Dockerfile}
which provisions them automatically.

\section{Installation}

The standard install path uses the project Makefile:

\begin{verbatim}
git clone https://github.com/IvanWeissVanDerPol/satellite-paraguay
cd satellite-paraguay
make install          # uv sync --all-extras + pip install -e .
make bootstrap        # install + data catalog + sample download
\end{verbatim}

The \texttt{bootstrap} target creates the \texttt{.venv/} virtual
environment, installs all Python dependencies pinned in
\texttt{uv.lock}, and downloads the sample Sentinel-2 + Hansen tiles
needed for the smoke tests. A full data bootstrap (real Hansen
v1.11 coverage of Paraguay + Sentinel-2 batch download) is the
\texttt{make data-sentinel} + \texttt{make data-mapbiomas} +
\texttt{make data-hansen} target chain.

\section{Verification}

Three verification commands establish that the installed environment
matches the canonical state of the repository:

\begin{itemize}
    \item \texttt{bash scripts/run\_tests.sh} — full pytest suite
      (~759 tests, ~80\,s wall-clock). Catches regressions in
      \texttt{src/}, \texttt{tests/}, \texttt{scripts/}, and the
      per-paper \texttt{paper.tex} / \texttt{references.bib} files.
    \item \texttt{bash scripts/run\_tests\_fast.sh} — fast inner-loop
      subset (7 test files, ~5\,s). Recommended for post-edit
      sanity checks during development.
    \item \texttt{bash scripts/defense\_check.py} — eight-check defense
      gate that bundles citation resolution, claim integrity,
      LaTeX syntax, ethics gates, and CrossRef DOI verification.
      This is the canonical pre-submission command.
\end{itemize}

All three commands are designed to exit non-zero on any failure; the
\texttt{run\_tests.sh} wrapper uses \texttt{set -euo pipefail} so a
single broken test blocks the whole run.

\section{Data Version Control}

Datasets larger than a few megabytes are tracked via DVC
(\texttt{dvc.yaml}) rather than committed to git. The canonical data
storage layer is the companion repository
\href{https://github.com/IvanWeissVanDerPol/paraguay-geodata-vlm}{\texttt{paraguay-geodata-vlm}}
(local path \texttt{/opt/data/thesis-active/}), which contains the
9-dataset manifest, OSM Paraguay downloader, and SAM +
GroundingDINO annotation pipeline. Per-paper \texttt{DATA\_MANIFEST.md}
files record which datasets are required for each paper; the
\texttt{scripts/validate\_data.py} worker audits whether the local
files actually exist or are missing/synthetic placeholders.

\section{Containerized Reproducibility}

For environments without \texttt{uv} or with restricted network
access, the bundled \texttt{Dockerfile} provides a fully
self-contained build:

\begin{verbatim}
docker build -t satellite-paraguay .
docker run -it --rm \
    -v $(pwd)/data:/app/data \
    satellite-paraguay bash scripts/run_tests.sh
\end{verbatim}

The image is multi-stage (Python 3.10-slim + GDAL system deps +
\texttt{uv}-installed Python deps) and produces a deterministic
environment that matches the \texttt{uv.lock} lock file. A
\texttt{docker-compose.production.yml} file is included for
deploying the FastAPI + Streamlit demo to a remote host.

\section{Continuous Integration}

The repository includes a GitHub Actions workflow
(\texttt{.github/workflows/ci.yml}) that runs the fast pytest subset
+ lint chain (black + flake8 + isort + mypy strict) on every push
and PR. A separate \texttt{latex.yml} workflow compiles each of the
six \texttt{paper.tex} files with TeX Live, ensuring that no paper
silently loses its ability to compile. The CI workflow SHAs were
pinned to v4.4.0 in commit \texttt{0b66057} after the v7.0.1
references were force-pushed upstream.
.
% Wordcount target: ~500 words of substantive prose + commands.

\chapter{Reproducibility Setup}
\label{ch:appendix:setup}

This appendix documents the complete reproducibility chain for the
six-paper thesis substrate: how to clone the repository, install the
Python dependencies, verify the installation, and run the full
verification suite. Every result in the thesis body is reproducible
from the public GitHub release by following the steps below; no
proprietary data or hidden state is required.

\section{Prerequisites}

The repository targets Python 3.10+ and depends on the
\texttt{uv} package manager for reproducible environment installation.
System-level dependencies (GDAL, GEOS, PROJ) are required for the
\texttt{rasterio} and \texttt{geopandas} imports used in the data
loaders; on Debian/Ubuntu these can be installed via the
\texttt{libgdal-dev} / \texttt{libgeos-dev} / \texttt{libproj-dev}
packages, or via the bundled multi-stage \texttt{Dockerfile}
which provisions them automatically.

\section{Installation}

The standard install path uses the project Makefile:

\begin{verbatim}
git clone https://github.com/IvanWeissVanDerPol/satellite-paraguay
cd satellite-paraguay
make install          # uv sync --all-extras + pip install -e .
make bootstrap        # install + data catalog + sample download
\end{verbatim}

The \texttt{bootstrap} target creates the \texttt{.venv/} virtual
environment, installs all Python dependencies pinned in
\texttt{uv.lock}, and downloads the sample Sentinel-2 + Hansen tiles
needed for the smoke tests. A full data bootstrap (real Hansen
v1.11 coverage of Paraguay + Sentinel-2 batch download) is the
\texttt{make data-sentinel} + \texttt{make data-mapbiomas} +
\texttt{make data-hansen} target chain.

\section{Verification}

Three verification commands establish that the installed environment
matches the canonical state of the repository:

\begin{itemize}
    \item \texttt{bash scripts/run\_tests.sh} — full pytest suite
      (~759 tests, ~80\,s wall-clock). Catches regressions in
      \texttt{src/}, \texttt{tests/}, \texttt{scripts/}, and the
      per-paper \texttt{paper.tex} / \texttt{references.bib} files.
    \item \texttt{bash scripts/run\_tests\_fast.sh} — fast inner-loop
      subset (7 test files, ~5\,s). Recommended for post-edit
      sanity checks during development.
    \item \texttt{bash scripts/defense\_check.py} — eight-check defense
      gate that bundles citation resolution, claim integrity,
      LaTeX syntax, ethics gates, and CrossRef DOI verification.
      This is the canonical pre-submission command.
\end{itemize}

All three commands are designed to exit non-zero on any failure; the
\texttt{run\_tests.sh} wrapper uses \texttt{set -euo pipefail} so a
single broken test blocks the whole run.

\section{Data Version Control}

Datasets larger than a few megabytes are tracked via DVC
(\texttt{dvc.yaml}) rather than committed to git. The canonical data
storage layer is the companion repository
\href{https://github.com/IvanWeissVanDerPol/paraguay-geodata-vlm}{\texttt{paraguay-geodata-vlm}}
(local path \texttt{/opt/data/thesis-active/}), which contains the
9-dataset manifest, OSM Paraguay downloader, and SAM +
GroundingDINO annotation pipeline. Per-paper \texttt{DATA\_MANIFEST.md}
files record which datasets are required for each paper; the
\texttt{scripts/validate\_data.py} worker audits whether the local
files actually exist or are missing/synthetic placeholders.

\section{Containerized Reproducibility}

For environments without \texttt{uv} or with restricted network
access, the bundled \texttt{Dockerfile} provides a fully
self-contained build:

\begin{verbatim}
docker build -t satellite-paraguay .
docker run -it --rm \
    -v $(pwd)/data:/app/data \
    satellite-paraguay bash scripts/run_tests.sh
\end{verbatim}

The image is multi-stage (Python 3.10-slim + GDAL system deps +
\texttt{uv}-installed Python deps) and produces a deterministic
environment that matches the \texttt{uv.lock} lock file. A
\texttt{docker-compose.production.yml} file is included for
deploying the FastAPI + Streamlit demo to a remote host.

\section{Continuous Integration}

The repository includes a GitHub Actions workflow
(\texttt{.github/workflows/ci.yml}) that runs the fast pytest subset
+ lint chain (black + flake8 + isort + mypy strict) on every push
and PR. A separate \texttt{latex.yml} workflow compiles each of the
six \texttt{paper.tex} files with TeX Live, ensuring that no paper
silently loses its ability to compile. The CI workflow SHAs were
pinned to v4.4.0 in commit \texttt{0b66057} after the v7.0.1
references were force-pushed upstream.
